% ==============================================================================
% APÊNDICE A — Prompts Completos Utilizados no Experimento
% ==============================================================================

\chapter{Prompts Completos Utilizados no Experimento}
\label{ap:prompts}

Este apêndice apresenta os prompts completos implementados no módulo de engenharia de prompts do agente (\texttt{prompt\_engine.py}), utilizados no experimento controlado (OE6), incluindo o system prompt e os templates de task prompts para cada tipo de mudança.

\section{System Prompt}

O system prompt define o papel do LLM e estabelece regras rígidas de comportamento, com ênfase explícita na preservação do estilo de aspas original (motivado por uma versão preliminar do experimento, em que o LLM ocasionalmente adicionava aspas simples desnecessárias ao redor de nomes de tabela):

\begin{lstlisting}[language=Python, caption=System Prompt, label=lst:system_prompt]
SYSTEM_PROMPT = """Você é um especialista em Power BI e DAX (Data Analysis Expressions).
Sua tarefa é refatorar expressões DAX preservando exatamente a mesma lógica de negócio.

REGRAS IMPORTANTES:
1. NUNCA altere a lógica de negócio - apenas aplique as mudanças estruturais solicitadas
2. Mantenha EXATAMENTE a formatação, indentação e estilo do original
3. Preserve todos os comentários existentes (linhas com //)
4. Se houver ambiguidade, escolha a interpretação mais segura
5. Retorne APENAS o código DAX refatorado, sem explicações adicionais
6. Se não for possível refatorar com segurança, retorne "ERRO: <motivo>"

REGRAS DE ESTILO OBRIGATÓRIAS:
- Use SEMPRE o mesmo formato de referência que está no código original
- Se o original usa  dGeral[Coluna]  -> mantenha  dGeral[Coluna]  (SEM aspas simples)
- Se o original usa  'dGeral'[Coluna] -> mantenha  'dGeral'[Coluna]  (COM aspas simples)
- NUNCA adicione aspas simples ao redor de nomes de tabela que nao tinham aspas no original
- NUNCA remova aspas simples de nomes de tabela que tinham aspas no original

FORMATO DE RESPOSTA:
Retorne apenas o codigo DAX entre tags <dax></dax>:
<dax>
-- codigo DAX aqui
</dax>"""
\end{lstlisting}

\section{Template: Renomeio de Coluna}

Template utilizado nos 8 casos de renomeio de coluna (D01--D08):

\begin{lstlisting}[language=Python, caption=Template de Renomeio de Coluna, label=lst:prompt_column]
TAREFA: Refatorar a expressao DAX abaixo para refletir o renomeio de uma coluna.

MUDANCA:
- Tabela: {table_name}
- Coluna antiga: [{old_name}]
- Coluna nova: [{new_name}]

INSTRUCAO CRITICA DE ESTILO:
Substitua APENAS as referencias a coluna renomeada.
Mantenha EXATAMENTE o mesmo formato de referencia de tabela que esta no codigo original.
Se o original usa  {table_name}[{old_name}]  (sem aspas), use  {table_name}[{new_name}]  (sem aspas).
NAO adicione aspas simples ao redor do nome da tabela se elas nao estavam no original.

EXEMPLO CORRETO:
  Original:  {table_name}[{old_name}]
  Resultado: {table_name}[{new_name}]

EXEMPLO ERRADO (nao faca isso):
  Original:  {table_name}[{old_name}]
  Resultado: '{table_name}'[{new_name}]  <- ERRADO, adicionou aspas desnecessarias

EXPRESSAO ORIGINAL:
```dax
{expression}
```

CONTEXTO ADICIONAL:
- Objeto: {object_name} ({object_type})
- Descricao: {description}

Retorne a expressao refatorada preservando o estilo original.
\end{lstlisting}

\section{Template: Renomeio de Tabela}

Template utilizado nos 4 casos de renomeio de tabela (RT01--RT04):

\begin{lstlisting}[language=Python, caption=Template de Renomeio de Tabela, label=lst:prompt_table]
TAREFA: Refatorar a expressao DAX abaixo para refletir o renomeio de uma tabela.

MUDANCA:
- Tabela antiga: '{old_table_name}'
- Tabela nova: '{new_table_name}'

SUBSTITUICOES NECESSARIAS:
- Todas as referencias a '{old_table_name}' devem ser substituidas por '{new_table_name}'
- Manter a mesma estrutura de aspas (se usava aspas simples, manter)

EXPRESSAO ORIGINAL:
```dax
{expression}
```

CONTEXTO:
- Objeto: {object_name} ({object_type})

Retorne a expressao refatorada.
\end{lstlisting}

\section{Template: Renomeio de Medida}

Template utilizado nos 8 casos de renomeio de medida (RM01--RM08):

\begin{lstlisting}[language=Python, caption=Template de Renomeio de Medida, label=lst:prompt_measure]
TAREFA: Refatorar a expressao DAX abaixo para refletir o renomeio de uma medida.

MUDANCA:
- Medida antiga: [{old_measure_name}]
- Medida nova: [{new_measure_name}]

EXPRESSAO ORIGINAL:
```dax
{expression}
```

CONTEXTO:
- Objeto: {object_name} ({object_type})

Retorne a expressao refatorada.
\end{lstlisting}

\section{Exemplo de Prompt Completo (Caso D01)}

Abaixo, um exemplo real de prompt enviado ao LLM durante o experimento, referente ao Caso D01 (renomeio da coluna \texttt{Plataforma}, impactando a medida \texttt{Mediana Segundos}):

\begin{lstlisting}[language=Python, caption=Exemplo de Prompt Completo (Caso D01), label=lst:prompt_example]
TAREFA: Refatorar a expressao DAX abaixo para refletir o renomeio de uma coluna.

MUDANCA:
- Tabela: dGeral
- Coluna antiga: [Plataforma]
- Coluna nova: [PlataformaUpdated]

INSTRUCAO CRITICA DE ESTILO:
Substitua APENAS as referencias a coluna renomeada.
Mantenha EXATAMENTE o mesmo formato de referencia de tabela que esta no codigo original.

EXPRESSAO ORIGINAL:
```dax
SWITCH(
 TRUE(),
 SELECTEDVALUE(dGeral[Plataforma]) = "Jenkins",
 ...
)
```

CONTEXTO ADICIONAL:
- Objeto: Mediana Segundos (measure)

Retorne a expressao refatorada preservando o estilo original.
\end{lstlisting}

\textbf{Resposta do LLM (Groq Llama 3.3 70B)}: expressão idêntica ao gabarito, com todas as 4 ocorrências de \texttt{dGeral[Plataforma]} substituídas por \texttt{dGeral[PlataformaUpdated]}, preservando integralmente a estrutura \texttt{SWITCH}/\texttt{CALCULATE}/\texttt{MEDIANX} original.

\textbf{Status}: Correto (Fase 1 e Fase 2), tempo de resposta de 1,84 segundos.

\section{Observações sobre o Prompt}

O system prompt inclui uma instrução explícita de estilo (preservação de aspas simples) e um exemplo de erro comum (``EXEMPLO ERRADO''), estratégia validada pela literatura como eficaz para reduzir erros de formatação em geração de código assistida por LLM (Seção~2.6). No experimento final, com 32 casos de teste, essa instrução contribuiu para uma taxa de erro de geração de 0\% nos 20 casos de renomeio.